\section{Background and related work}
\label{sec:related}

\Cref{tab:families} groups prior work into five lines of research and lists the properties that matter in our
setting. No line covers all of them; \nova{} borrows one element from each.

\begin{table}[t]
\centering
\caption{Lines of prior work and the properties they typically address (\checkmark: usually addressed;
$\circ$: partly or in some works; --: usually not addressed).}
\label{tab:families}
\small
\resizebox{\textwidth}{!}{%
\begin{tabular}{@{}lcccccc@{}}
\toprule
Line of work & \makecell{Rejects\\unknown flows} & \makecell{Adapts\\the model} & \makecell{Uses analyst\\labels} &
\makecell{Models label\\delay} & \makecell{Validates\\updates} & \makecell{Multi-class\\attacks} \\
\midrule
Drift detection and rejection~\cite{jordaney2017transcend,barbero2022transcending,yang2021cade,han2023owad}
  & \checkmark & $\circ$ & $\circ$ & -- & -- & $\circ$ \\
Self-training adaptive NIDS~\cite{andresini2021insomnia,zhang2024aoc,zhang2025ssf,mousavipour2025shift}
  & -- & \checkmark & $\circ$ & -- & -- & -- \\
Active learning under drift~\cite{chen2023continuous,gupta2025netguard,prabhu2021clue}
  & -- & \checkmark & \checkmark & -- & -- & $\circ$ \\
Continual learning with replay~\cite{amalapuram2023augmented,amalapuram2024soul,chrysakis2020online}
  & $\circ$ & \checkmark & $\circ$ & -- & -- & $\circ$ \\
Open-set NIDS~\cite{zhang2020osr,corsini2023ood}
  & \checkmark & -- & -- & -- & -- & \checkmark \\
\midrule
\nova{} (this work) & \checkmark & \checkmark & \checkmark & \checkmark & \checkmark & \checkmark \\
\bottomrule
\end{tabular}}
\end{table}

\subsection{Evaluation methodology and datasets}
TESSERACT~\cite{pendlebury2019tesseract} showed that two biases inflate results of security classifiers:
\emph{spatial} bias, when the class ratio of the test set is unrealistic, and \emph{temporal} bias, when training
data postdate test data. It proposed time-consistent splits and the area-under-time (AUT) metric, which we adopt.
Arp et al.~\cite{arp2022dos} catalogue ten pitfalls of machine learning in security, among them data snooping,
inappropriate baselines and lab-only evaluation; our protocol is designed against them (\cref{sec:problem}).

CIC-IDS2017~\cite{sharafaldin2018} and CSE-CIC-IDS2018 are the most widely used flow datasets, and both contain
errors. Engelen et al.~\cite{engelen2021troubleshooting} found feature-extraction artefacts and mislabelled flows
in CIC-IDS2017. Lanvin et al.~\cite{lanvin2022errors} independently found unlabelled attack traffic and showed
that the corrections change reported detection results. Liu et al.~\cite{liu2022error} audited both CIC
datasets. UNSW-NB15~\cite{moustafa2015unsw} was captured on two days with a different tool chain. We use the
three datasets as published and discuss the consequences of their errors in \cref{sec:discussion}. Performance
estimates from chronological splits of these datasets differ from those of random
splits~\cite{nismara2026temporal}, and a recent sequence study on CIC-IDS2017 found that the protocol matters far
more than the model architecture~\cite{moczkodan2026transformers}.

\subsection{Drift detection and rejection}
Transcend~\cite{jordaney2017transcend} attaches conformal credibility scores to the predictions of any classifier
and rejects low-credibility ones as drifting; its successor~\cite{barbero2022transcending} makes the calibration
cheaper. CADE~\cite{yang2021cade} trains a contrastive autoencoder and flags samples whose distance to every class
centroid, scaled by the class's spread, is large; it also explains which features cause the drift. OWAD~\cite{han2023owad}
targets anomaly detectors trained on benign traffic: it calibrates their outputs, detects normality shift with a
permutation test, selects a few samples that explain the shift for labelling, and adapts with a regulariser
against forgetting. These works decide \emph{whether} the input distribution has moved. \nova{} uses a
rejection rule of the same kind, but conditions it on the predicted class (a single centroid cannot represent
multi-modal benign traffic, \cref{sec:res-ablation}) and couples it with labelling and deployment decisions.

\subsection{Adaptive network intrusion detection}
INSOMNIA~\cite{andresini2021insomnia} retrains a binary detector on drifting CIC-IDS2017 traffic. In every window,
the least-confident flows are labelled by a nearest-centroid estimator instead of an analyst, and both models are
fine-tuned. Because the estimator can only produce classes seen in training, genuinely new attacks are never
learned, and its errors are fed back into training. AOC-IDS~\cite{zhang2024aoc} fine-tunes a contrastive
autoencoder online on pseudo-labels from a Gaussian fit of similarity to a benign prototype. SSF~\cite{zhang2025ssf}
adds drift detection, a learned selection of old and new samples, a small labelled share and strategic forgetting.
Mousavipour et al.~\cite{mousavipour2025shift} detect shift at the distribution level, explain it with masks and
adapt with relational distillation. ENIDrift~\cite{wang2022enidrift} keeps an ensemble of sub-classifiers over
incremental packet embeddings and prunes unstable members. Most of these methods are binary, rely on their own
pseudo-labels, and were evaluated on random or day-wise splits without label delay. We re-implement INSOMNIA as a
representative of the self-training line.

\subsection{Active learning under drift}
Chen et al.~\cite{chen2023continuous} retrain an Android-malware classifier monthly and choose samples to label
with a pseudo-loss computed from nearest training neighbours in a hierarchically contrastive embedding; they show
that similarity-based selection beats plain uncertainty under drift. NetGuard~\cite{gupta2025netguard} adapts a
NIDS from CIC-IDS2017 to CSE-CIC-IDS2018: it labels the flows with the highest density ratio between new and
labelled traffic (Gaussian mixtures on each) and synthesises rare attacks with a tabular generator.
CLUE~\cite{prabhu2021clue} clusters uncertainty-weighted embeddings and labels the sample nearest each centroid,
which yields uncertain and diverse picks. \nova{}'s query belongs to this family, but it weights clusters by an
open-world novelty score, keeps a uniform share for benign drift, and runs every window under a label delay. We
re-implement uncertainty sampling and NetGuard's selection.

\subsection{Continual learning with replay}
Continual learning keeps a memory of past samples to avoid forgetting. Class-balancing reservoir sampling
(CBRS)~\cite{chrysakis2020online} lets rare classes survive in a fixed memory. Amalapuram et
al.~\cite{amalapuram2023augmented} extend it for intrusion detection and make interference-aware retrieval
cheaper; SOUL~\cite{amalapuram2024soul} adds semi-supervised label generation for unseen tasks, keeping a
pseudo-label only if it agrees with a similarity vote over the memory. We re-implement CBRS replay, which turned
out to be the strongest baseline, and \nova{}'s label propagation follows SOUL's agreement principle.

\subsection{Open-set recognition and out-of-distribution detection}
Open-CNN~\cite{zhang2020osr} fits a Weibull distribution to the activations of known classes and estimates an
unknown-class probability. Corsini and Yang~\cite{corsini2023ood} compare out-of-distribution detectors on the CIC
datasets and find the distance to the $k$-th nearest training embedding the strongest single detector, in line
with deep $k$-NN results in vision~\cite{sun2022knn}. These detectors were evaluated on static splits.
\nova{} uses a class-conditional $k$-NN detector and recalibrates it whenever a new class is learned; we show
that the recalibration is what keeps learned classes from being rejected after drift (\cref{sec:res-main}).
Few-shot open-set recognition~\cite{liu2020peeler} learns embeddings by episodic meta-learning; we test whether
such an embedding would help in \cref{sec:res-meta}.
